A variational characterization of the quantum fidelity between two density operators $\rho$ and $\sigma$ is given by Alberti's theorem:
$$
F(\rho,\sigma)^2
=
\inf_{P>0}
\Tr(\rho P)\Tr(\sigma P^{-1}),
$$
where the infimum is taken over all invertible positive-definite operators $P$.

Assume that $\rho$ and $\sigma$ commute. Hence, there exists an orthonormal basis $\{\ket{i}\}$ such that
$$
\rho=\sum_i \lambda_i \ket{i}\bra{i},
\qquad
\sigma=\sum_i \mu_i \ket{i}\bra{i},
$$
where $\lambda_i,\mu_i\geq 0$ and
$$
\sum_i\lambda_i=\sum_i\mu_i=1.
$$

Restrict $P$ to be diagonal in the common eigenbasis,
$$
P=\sum_i p_i \ket{i}\bra{i},
\qquad
p_i>0.
$$
Using the Cauchy--Schwarz inequality, solve the resulting optimization problem and prove that
$$
\inf_{\{p_i>0\}}
\left(
\sum_i \lambda_i p_i
\right)
\left(
\sum_i \mu_i p_i^{-1}
\right)
=
\left(
\sum_i \sqrt{\lambda_i\mu_i}
\right)^2.
$$
Conclude that, in the commuting case,
$$
F(\rho,\sigma)
=
\sum_i \sqrt{\lambda_i\mu_i},
$$
which is the classical fidelity, also known as the Bhattacharyya coefficient.
